%!TEX root = ../dissertation.tex
% !TeX spellcheck = en_GB 

\chapter{State-of-the-art algorithms}\label{chapt:2}

\section{Random Search}\label{sect:Nesterov}

The first relevant algorithm proposed in the literature to approach the problem (\ref{problem}), is the algorithm from \cite{nesterov2017random}, \textit{Random Search}, with the aim to simulate
the working of the gradient descent but translated in the zeroth-order optimization.

In order to overcome the limitations posed by the impossibility of computing the gradient, the idea of Nesterov, and the
literature that followed, was to replace the objective function with a new function which keeps most of the
properties of the original but makes the computation of the gradient possible. This strategy is the so called Gaussian Smoothing:

\begin{defn}[Gaussian Smoothing]\label{GaussianSmoothingDefn}

	Let $\sigma > 0$ the \textit{smoothing parameter}, then we denote by $F_\sigma (x)$ the Gaussian Smoothing of $F$ with radius $\sigma$, i.e.
	
	\begin{equation}
		F_{\sigma}(x)=\frac{1}{\pi^{d / 2}} \int_{\mathbb{R}^{d}} F(x+\sigma \xi) e^{-\|\xi\|_{2}^{2}} \mathrm{~d} \xi=\mathbb{E}_{\xi \sim \mathcal{N}\left(0, \mathbb{I}_{d}\right)}[F(x+\sigma \xi)]
	\end{equation}
	
\end{defn}

$F_\sigma$ keeps all the features of $F$, such as convexity and the Lipschitz constant, while being always differentiable, even when $F$ is not. This is relevant because, given that $\norm{F - F_\sigma}$ can be bounded by the Lipschitz constant, (\ref{problem}) can be replaced with

\begin{equation}\label{SmoothedProblem}
    \min _{x \in \mathbb{R}^{d}} F_\sigma (x)  
\end{equation}

Now, the gradient of the smoothed $F$ is:

\begin{equation}
    \nabla F_{\sigma}(x)=\frac{2}{\sigma \pi^{d / 2}} \int_{\mathbb{R}^{d}} \xi F(x+\sigma \xi) e^{-\|\xi\|_{2}^{2}} \mathrm{~d} \xi=\frac{2}{\sigma} \mathbb{E}_{\xi \sim \mathcal{N}\left(0, \mathbb{I}_{d}\right)}[\xi F(x+\sigma \xi)]
\end{equation}

This is still a difficult object to compute, but it can be estimated easily. The estimator usually exploited to estimate this gradient is obtained via Monte Carlo sampling, with a central finite difference approach, which is used to reduce the variance of the estimation; namely, given $M \in \mathbb{N}_+$ the number of directions and $\left\{\xi_{j}\right\}_{j=1}^{M} \stackrel{\mathrm{iid}}{\sim} \mathcal{N}(0,1)$ the directions, then we have the following:

\begin{defn}[Gradient of Gaussian Smoothed function]\label{GradientGaussianSmoothingDefn}

	\begin{equation}
	    \nabla F_{\sigma}(x) \approx \frac{1}{2\sigma M} \sum_{j=1}^{M} \xi_{j}\left(F\left(x+\sigma \xi_{j}\right)-F\left(x-\sigma \xi_{j}\right)\right)
	\end{equation}

\end{defn}

Actually, Nesterov in its work first proposed and studied the theoretical properties of its algorithm with $M=1$;
following literature increased it.

\begin{algorithm}[h]
\caption{Random Search}\label{NesterovAlgorithm}

\hspace*{\algorithmicindent} \textbf{Input:} function $f$, state $x_0$, smoothing parameter $\sigma$, stepsize $\lambda$
\begin{algorithmic}[1]

\For{$t=0,1,2,\dots$}
    \State Sample $\xi \sim \mathcal{N}(0,I_n)$
    \State Compute returns $F_{+} = F(x_t + \sigma\xi)$, $F_{-} = F(x_t - \sigma\xi)$
    \State Set $x_{t+1} = x_t - \frac{\lambda}{2\sigma} (F_{t_+} - F_{t_-})\xi$
\EndFor

\end{algorithmic}
\end{algorithm}

The main theoretical result proved by Nesterov for smooth convex functions is the following. 

At first, note that this method generates a random sequence $ \left\{x_{k}\right\}_{k \geq 0} . $ Denote by

\begin{equation}
	\mathcal{U}_{k}=\left(\xi_{0}, \ldots, \xi_{k}\right)
\end{equation}

a random vector composed by i.i.d. variables $ \left\{u_{k}\right\}_{k \geq 0} $ related to each iteration of the scheme. Denote $ \phi_{0}=F\left(x_{0}\right) $, and $ \phi_{k} \stackrel{\text { def }}{=} \mathbb{E}_{\mathcal{U}_{k-1}}\left[ F\left(x_{k}\right)\right],\quad k \geq 1 $. Then it holds:

\begin{thm}
	Let $ F $ have Lipschitz Continuos Gradient, and sequence $ \left\{x_{k}\right\}_{k \geq 0} $ be generated by \ref{NesterovAlgorithm} with
	
	\begin{equation}
		\lambda=\frac{1}{4(d+4) L}
	\end{equation}

	Then, for any $ N \geq 0 $, we have
	
	\begin{equation}
		\frac{1}{N+1} \sum_{k=0}^{N}\left(\phi_{k}-f^{*}\right) \leq \frac{4(d+4) L\left\|x_{0}-x^{*}\right\|^{2}}{N+1}+\frac{9 \sigma^{2}(d+4)^{2} L}{25}
	\end{equation}

	Let function $F$ be strongly convex. Denote $ \delta_{\sigma}=\frac{18 \sigma^{2}(d+4)^{2}}{25 \tau} L . $ Then
	
	\begin{equation}
		\phi_{N}-F^{*} \leq \frac{1}{2} L\left[\delta_{\sigma}+\left(1-\frac{\tau}{8(d+4) L}\right)^{N}\left(\left\|x_{0}-x^{*}\right\|^{2}-\delta_{\sigma}\right)\right]
	\end{equation}

\end{thm}

%Add some brief comments on the theoretical result.
Considering the strongly convex case, the smoothing parameter $\mu$ has to be chosen such that it satisfies the inequality $\frac{1}{2}L\delta_\mu$. In this case, the number of iterations needed by the method to satisfy $\phi_{N}-F^{*} \leq \varepsilon$, is of the order of $\mathcal{O}(d\log\frac{1}{\varepsilon})$. It is interesting to note that the original gradient method requires $\mathcal{O}(\log\frac{1}{\varepsilon})$ iterations, which shows that this method, and all the other methods presented in this work, suffer from the curse of dimensionality. This problem can not be avoided: in fact using $M=d$ directions to estimate the gradient would lead to $\mathcal{O}(\log\frac{1}{\varepsilon})$ iterations, at the cost of increasing the computational time for each iteration by a factor of $d$.

%-------------------------------------------------------------------------------

\section{ES}\label{ESAlgorithm}

The following algorithm, from \cite{Salimans2017}, is essentially an analogous of Nesterov's Random Search which has its main focus
on the increment on the number $M$ of directions used during the estimation of the gradient.

\begin{algorithm}[h]
\caption{ES algorithm}\label{gd/es}

\hspace*{\algorithmicindent} \textbf{Input:} function $f$, state $x_0$, smoothing parameter $\sigma_0$, stepsize $\lambda$
\begin{algorithmic}[1]

\For{$t=0,1,2,\dots$}
    \State Sample $\xi_1,\dots,\xi_n \sim \mathcal{N}(0,I_n)$
    \State Compute returns $F_{i_+} = F(x_t + \sigma\xi_i)$, $F_{i_-} = F(x_t - \sigma\xi_i)$, for $i=1,\dots,n$
    \State Set $x_{i+1} = x_i - \frac{\lambda}{2n\sigma}\sum_{i=1}^n (F_{i_+} - F_{i_-})\xi_i$
\EndFor

\end{algorithmic}
\end{algorithm}

The authors focused on the use of the algorithm to solve RL problems, as in \ref{problem}, and the possibility of parallelizing the execution of the algorithm which showed the
strength of the algorithm to train AI's for this aim. This is proven to be able to led to an incredible advantage against other RL algorithms.

%--------------------------------------------------------------------------

\section{ASEBO}\label{sect:ASEBO}

Inspired by the classical RL literature, in \cite{Choromanski2019}, the authors tried to embed the idea of the
exploration-exploitation tradeoff into the framework given by Algorithm \ref{NesterovAlgorithm} and \ref{gd/es}. To
achieve this they introduced the idea of active subspaces, originally presented in \cite{constantine2015active}, and they devised an adaptive
method that tries to balance between exploration and exploitation. In their work they identified exploration as
the choice of directions independent of previous gradients directions, while exploitation was associated to choice of
random directions dependent of previous gradients.

\begin{algorithm}[h]
\caption{ASEBO algorithm}\label{ASEBO}

\hspace*{\algorithmicindent} \textbf{Input:} function $F$, state $x_0$, smoothing parameter $\sigma_0$, stepsize $\lambda$, decay rate $\nu$, termination tolerance $\varepsilon$, number of iterations of full-sampling $l$, total iterations \textit{maxiter}\\
\hspace*{\algorithmicindent} \textbf{Output:} point of minimum $\Tilde{x}$
\begin{algorithmic}[1]

\State Set $\Tilde{x} = x_0$, $\sigma = \sigma_0$ and let $\Sigma$ be a random orthonormal basis
\State Initialize an archive $\mathcal{G}$ of maximum capacity $k$
\For{$t=0$ to \textit{maxiter}}
    \If{$t<l$}
        \State Let $n_t = d$. Sample $\xi_1,\dots,\xi_{n_t}$ from $\mathcal{N}(0,I_d)$
    \Else
        \State \multiline{Take top $ r $ eigenvalues $ \lambda_{i} $ of $ \mathrm{Cov}_{t} $, where $ r $ is smallest such that: $ \sum_{i=1}^{r} \dot{\lambda}_{i} \geq \epsilon \sum_{i=1}^{d} \lambda_{i} $ using its SVD as described in text and take $ n_{t}=r $.}
        \State \multiline{Take the corresponding eigenvectors $ \mathbf{u}_{1}, \ldots, \mathbf{u}_{r} \in \mathbb{R}^{d} $ and let $ \mathbf{U} \in \mathbb{R}^{d \times r} $ be obtained by stacking them together. Let $ \mathrm{U}^{\mathrm{act}} \in \mathbb{R}^{d \times r} $ be obtained from stacking together some orthonormal basis of $ \mathcal{L}_{\text {active }}^{\mathrm{ES}} \stackrel{\text { def }}{=} \operatorname{span}\left\{\mathbf{u}_{1}, \ldots, \mathbf{u}_{r}\right\} . $ Let $ \mathbf{U}^{\perp} \in \mathbb{R}^{d \times(d-r)} $ be obtained from stacking together some orthonormal basis of the orthogonal complement $ \mathcal{L}_{\text {active }}^{\mathrm{ES}, \perp} $ of $ \mathcal{L}_{\text {active }}^{\text {ES }} $}
        \State \multiline{Sample $ n_{t} $ vectors $ \mathbf{g}_{1}, \ldots, \mathbf{g}_{n_{t}} $ as follows: with probability $ 1-p^{t} $ from $ \mathcal{N}(0, \mathbf{U}^{\perp}(\mathbf{U}^{\perp})^{\top}) $
        	and with probability $ p^{t} $ from $ \mathcal{N}(0, \mathbf{U}^{\text {act }}(\mathbf{U}^{\text {act }})^{\top}) $.}
        \State Renormalize $ \mathbf{g}_{1}, \ldots, \mathbf{g}_{n_{t}} $ such that marginal distributions $ \left\|\mathbf{g}_{i}\right\|_{2} $ are $ \chi(d) $.
    \EndIf
    \State Compute $ \widehat{\nabla}_{\mathrm{MC}}^{\mathrm{AT}} F\left(x_{t}\right) $ as: $ \widehat{\nabla}_{\mathrm{MC}}^{\mathrm{AT}} F\left(x_{t}\right)=\frac{1}{2 n_{t} \sigma} \sum_{j=1}^{n_{t}}\left(F\left(x_{t}+\mathbf{g}_{j}\right)-F\left(x_{t}-\mathbf{g}_{j}\right)\right) \mathrm{g}_{j} $
    \State Set $ \operatorname{Cov}_{t+1}=\lambda \operatorname{Cov}_{t}+(1-\lambda) \Gamma $, where $ \Gamma=\widehat{\nabla}_{\mathrm{MC}}^{\mathrm{AT}} F_{\sigma}\left(x_{t}\right)\left(\widehat{\nabla}_{\mathrm{MC}}^{\mathrm{AT}} F_{\sigma}\left(x_{t}\right)\right)^{\top} . $
    \State Set $ p^{t+1}=p_{\text {opt }} $ for $ p_{\text {opt }} $ output by Algorithm 2 and: $ x_{t+1} = x_{t} - \eta \widehat{\nabla}_{\mathrm{MC}}^{\mathrm{AT}} F\left(x_{t}\right) $.
    
\EndFor

\end{algorithmic}
\end{algorithm}

\begin{defn}\label{SmoothThirdOrderDerivative}
	
	$ F $ has a $ \tau $-smooth third order derivative tensor with respect to $ \sigma>0 $, so that $ F(x+\sigma \mathbf{g})=F(x)+\sigma \nabla F(x)^{\top} \mathbf{g}+\frac{\sigma^{2}}{2} \mathbf{g}^{\top} H(x) \mathbf{g}+\frac{1}{6} \sigma^{3} F^{\prime \prime \prime}(x)[\mathbf{v}, \mathbf{v}, \mathbf{v}] $ for some $ \mathbf{v} \in \mathbb{R}^{d} $
	$ \left(\|\mathbf{v}\|_{2} \leq\|\mathbf{g}\|_{2}\right) $ satisfying $ \mid F^{\prime \prime \prime}(x)[\mathbf{v}, \mathbf{v}, \mathbf{v}] \leq \tau\|\mathbf{v}\|_{2}^{3} \leq \tau\|\mathbf{g}\|_{2}^{3} $

\end{defn}

The authors proved the following bounds:

\begin{thm}

	If $ F $ has Lipschitz Continuos Gradient and $ \tau $-smooth third order derivative tensor, the estimators $ \widehat{\nabla}_{\text {MC }, k=1}^{\text {AT,base }} F_{\sigma}(x) $ and $ \widehat{\nabla}_{\text {MC, } k=1}^{\text {AT,asebo }} F_{\sigma}(x) $ are close to the true gradient $ \nabla F(x) $, i.e.: 
	
	\begin{equation}
		\left\|\mathbb{E}_{\mathbf{g} \sim \mathcal{N}\left(0, \mathbf{I}_{d}\right)}\left[\widehat{\nabla}_{\mathrm{MC}, k=1}^{\mathrm{AT}, \text { base }} F_{\sigma}(x)\right]-\nabla F(x)\right\| \leq \epsilon
	\end{equation}

	and

	\begin{equation}
		\left\|\mathbb{E}_{\mathbf{g} \sim \widehat{P}}\left[\widehat{\nabla}_{\mathbf{M C}, k=1}^{\mathrm{AT}, \text { asebo }} F_{\sigma}(x)\right]-\nabla F(x)\right\| \leq \epsilon
	\end{equation}

\end{thm}

\begin{algorithm}[h]
\caption{ASEBO subroutine}\label{ASEBO subroutine}

\hspace*{\algorithmicindent} \textbf{Hyperparameters:} smoothing parameter $ \sigma $, horizon $ C $, learning rate $ \alpha $, probability regularizer $ \beta $, initial probability parameter $ q_{0}^{t} \in(0,1) $\\
\hspace*{\algorithmicindent} \textbf{Input:} subspaces: $ \mathcal{L}_{\text {active }}^{\mathrm{ES}}, \mathcal{L}_{\text {active }}^{\mathrm{ES}, \perp} $, function $ F $, vector $ x_{t} $\\
\hspace*{\algorithmicindent} \textbf{Output:} $p_C$
\begin{algorithmic}[1]

\For{$t=0$ to \textit{maxiter}}
    \State Compute $ p_{l-1}^{t}=(1-2 \beta) q_{l-1}^{t}+\beta $ and sample $ a_{l}^{t} \sim \operatorname{Ber}\left(p_{l}^{t}\right) $
    \State If $ a_{l}^{t}=1 $, sample $ \mathbf{g}_{l} \sim \mathcal{N}\left(0, \sigma \mathbf{I}_{\mathcal{L}_{\text {active }}^{\text {ES }}}\right) $, otherwise sample $ \mathbf{g}_{l} \sim \mathcal{N}\left(0, \sigma \mathbf{I}_{ \mathcal{L}_{\text {active }}^{\text{ES, }\perp}}\right) $
    \State Compute $ v_{l}=\frac{1}{2 \sigma}\left(F\left(x_{t}+\mathbf{g}_{l}\right)-F\left(x_{t}-\mathbf{g}_{l}\right)\right) $
    \State Set $ \mathbf{e}_{l}=(1-2 \beta)\left[\begin{array}{c}\left(-\frac{a_{l}^{t}\left(\operatorname{dim}\left(\mathcal{L}_{\text {ative }}^{\mathrm{ES}}\right)+2\right)}{\left(p_{l}^{t}\right)^{3}}\right) \\ \left(-\frac{\left(1-a_{l}^{t}\right)\left(\operatorname{dim}\left(\mathcal{L}_{\text {active }}\right)+2\right)}{\left(1-p_{l}^{t}\right)^{3}}\right)\end{array}\right] v_{l}^{2} $
    \State Set $ q_{l}^{t}=\frac{q_{l-1}^{t} \exp \left(-\alpha \mathbf{e}_{l}(1)\right)}{q_{l-1}^{t} \exp \left(-\alpha_{l}(1)\right)+\left(1-q_{l-1}^{t}\right) \exp \left(-\alpha \mathrm{e}_{l}(2)\right)} $
\EndFor

\end{algorithmic}
\end{algorithm}

\begin{thm}
	
	The following holds for $ s_{\mathbf{U}^{\text {act }}}=\left\|\left(\mathbf{U}^{\mathrm{act}}\right)^{\top} \nabla
	F(x)\right\|_{2}^{2} $ and $ s_{\mathbf{U}^{\perp}}=\left\|\left(\mathbf{U}^{\perp}\right)^{\top} \nabla
	F(x)\right\|_{2}^{2} :$ 
	
	\begin{itemize}

		\item The variance of $ \widehat{\nabla}_{\mathrm{MC}, k=1}^{\mathrm{AT} \text { ,asebo }} F_{\sigma}(x) $ is close to $ \Gamma $, i.e. $ \left|\operatorname{Var}\left[\widehat{\nabla}_{\mathrm{MC}, k=1}^{\mathrm{AT}, \text { asebo }} F_{\sigma}(x)\right]-\Gamma\right| \leq \epsilon $
		
		\item The choice of $ p^{t} $ that minimizes $ \Gamma $ satisfies $ p_{*}^{t}:=\frac{\sqrt{\left(s_{\text {Uact }}\right)\left(d_{\text {active }+2)}\right.}}{\sqrt{\left(s_{\text {Uact }}\right)\left(d_{\text {active }}+2\right)}+\sqrt{\left(s_{\mathrm{U}} \perp\right)\left(d_{\mathrm{U} \perp}+2\right)}} $ and the optimal variance  $Var _{\text {opt }} $ corresponding to $ p_{*}^{t} $ satisfies: $ \mid $  $Var _{\text {opt }}-\Delta \mid \leq \epsilon $ for 
		\begin{equation*}
			\Delta= \left[\sqrt{\left(s_{ \left.\mathbf{U}^{\text {act }}\right)\left(d_{\text {active }}+2\right)}\right.}+\sqrt{\left(s_{ \left.\mathbf{U}^{\perp}\right)\left(d_{\perp}+2\right)}\right.}\right]^{2}-\|\nabla F(x)\|^{2}
		\end{equation*}
	
		\item $ \mathrm{Var}_{\mathrm{opt}} \leq \operatorname{Var}\left[\widehat{\nabla}_{\mathrm{MC}, k=1}^{\mathrm{AT}, \text { base }} F_{\sigma}(x)\right]+\epsilon\\ -\underbrace{\left|\sqrt{\left(s_{\mathbf{U}^{\perp}}\right)\left(d_{\text {active }}+2\right)}-\sqrt{\left(s_{\mathbf{U}^{\mathrm{act}}}\right)\left(d_{\perp}+2\right)}\right|^{2}-2\|\nabla F(x)\|^{2}}_{\lambda} $.

	\end{itemize}

\end{thm}

%---------------------------------------------------------------------------

\section{SGES}\label{sect:SGES}

Following the work of \cite{Choromanski2019}, \cite{Liu2020a} tried to improve and simplify Algorithm \ref{ASEBO}. 

Instead of keeping
track of all the gradients via momentum-like update of a covariance matrix, a fixed-size buffer of gradients is kept
updated and used to compute a covariance matrix; the routine responsible for the choice of directions is also simplified
and more numerically stable. To do so, $k\ll d$ gradients are collected from the beginning, and then always the most recent $k$ gradients in the buffer are kept. Then at each iteration $t \geq k$, given $\mathbf{G}_t\in \mathbb{R}^{n\text{x} k}$, which is a matrix made with the last $k$ gradients, the algorithm generates the subspace $\mathcal{L}_\mathbf{G}$, using the vectors in the buffer. This subspace captures the structure of the optimization problem and, in the reinforcement learning optic, it captures the directions related to exploitation.
Then the algorithm samples the search directions from:

\begin{equation}\label{historical sampling}
	\left\{\begin{array}{ll}\xi \sim \mathcal{N}\left(0, \mathrm{G}^{\perp} \mathrm{G}\right) & \text { with probability } \alpha \\ \xi \sim \mathcal{N}(0, \mathrm{I}) & \text { with probability } 1-\alpha\end{array}\right. 
\end{equation}

$\alpha\in\left( 0,1 \right)$ represents the trade-off between exploitation, choosing a direction from $L_\mathbf{G}$, and exploration, from the whole space. The algorithm adapts $\alpha$ at each iteration such that the amount of exploitation and exploration varies depending on how the algorithm is performing. To do that the following are computed:

\begin{equation}
	\widehat{r}_{\mathbf{G}}=\frac{1}{M} \sum_{i=1}^{M} \min_{m=1,\dots,m_i} \left\{ F\left( x+\sigma \xi_i \right) \right\} 
\end{equation}

\begin{equation}
	\widehat{r}_{\mathbf{G}}^{\perp}=\frac{1}{d-M} \sum_{i=M+1}^{d} \min_{m=1,\dots,m_i} \left\{ F\left( x+\sigma \xi_i \right) \right\} 
\end{equation}

Then, at iteration $t$, $\alpha$ is chosen as:

\begin{equation}\label{adapt alpha}
	\alpha_{t}=\left\{\begin{array}{ll}\min \left\{\delta \alpha_{t-1}, \kappa_{1}\right\} \text { if } \widehat{r}_{\mathbf{G}} \text { does not exist or } \widehat{r}_{\mathbf{G}}<\widehat{r}_{\mathbf{G}}^{\perp} \\ \max \left\{\frac{1}{\delta} \alpha_{t-1}, \kappa_{2}\right\} \text { if } \widehat{r}_{\mathbf{G}}^{\perp} \text { does not exist or } \widehat{r}_{\mathbf{G}} \geq \widehat{r}_{\mathbf{G}}^{\perp}\end{array}\right. 
\end{equation}

with $\delta > 1$ a scaling factor, $k_1$ and $k_2$ are upper and lower bounds of $\alpha$. Intuitively, if $\widehat{r}_{\mathrm{G}}<\widehat{r}_{\mathrm{G}}^{\perp} $ or $\widehat{r}_{\mathrm{G}}$ does not exist, than the algorithm increases $\alpha$, otherwise it is decreased.

\begin{algorithm}[h]
\caption{SGES algorithm}\label{SGES}

\hspace*{\algorithmicindent} \textbf{Input:} function $F$, state $x_0$, smoothing parameter $\sigma_0$, learning rate $\lambda$, hyper-parameters $k$, total iterations \textit{maxiter}, warmup iteration $T_w \geq k$\\
\hspace*{\algorithmicindent} \textbf{Output:} point of minimum $\Tilde{x}$
\begin{algorithmic}[1]

\State Initialize an archive $\mathcal{G}$ of maximum capacity $k$
\For{$t=0$ to \textit{maxiter}}
    \If{$t<T_w$}
        \State Sample search directions $\xi_1,\dots,\xi_n$ from $\mathcal{N}(0,I_n)$
    \Else
        \State Obtain gradient matrix $\mathbf{G}\in\mathbb{R}^{nxk}$ from $\mathcal{G}$
        \State Generate subspaces $L_{\mathbf{G}}$ and $L_{\mathbf{G}}^{\perp}$
        \State Sample search directions $\varepsilon_1,\dots,\varepsilon_n$ from (\ref{historical sampling})
        \State Normalize these search directions
    \EndIf
    \State Compute gradient estimate $ \nabla F_{\sigma}\left(x_{i} \right) $ by (\ref{GradientGaussianSmoothingDefn})
    \State Update the parameters $ x_{i+1}=x_{i}-\lambda \nabla F_{\sigma}\left(x_{i}\right) $
    \State Add the gradient estimate $ \nabla F_{\sigma}\left(x_{i} \right) $ to $\mathcal{G}$
    \If{$i < T_w$}
        \State Adaptively adjust $\alpha$ as in (\ref{adapt alpha})
    \EndIf
\EndFor

\end{algorithmic}
\end{algorithm}

Their main theoretical result is:

\begin{thm}
	
	Assuming the objective function has Lipschitz Continous Gradient, \ref{SmoothThirdOrderDerivative} and $ \sigma<\frac{1}{35} \sqrt{\frac{\varepsilon \min \{\alpha, 1-\alpha\}}{\tau n^{3} \max \{L, 1\}}} $ for some precision
	parameter $ \varepsilon>0 $, it holds

	\begin{equation}
		\left\|\mathbb{E}_{\boldsymbol{\epsilon} \sim \mathcal{P}}\left[\hat{g}_{\text {sges }}\right]-\nabla f(x)\right\| \leq \varepsilon \\
	\left|\operatorname{Var}\left[\hat{g}_{\text {sges }}\right]-\Omega\right| \leq \varepsilon
	\end{equation}
	
	where
	
	\begin{equation}
		\Omega=\frac{k+2}{\alpha} \cdot\left\|\mathbf{U}^{\top} \nabla F(x)\right\|^{2}-\|\nabla F(x)\|^{2} + \frac{n-k+2}{1-\alpha} \cdot\left\|\left(\mathbf{U}^{\perp}\right)^{\top} \nabla F(x)\right\|^{2}
	\end{equation}

\end{thm}

%-------------------------------------------------------------------------

\section{ASGF}\label{sect:ASGF} 

This algorithm from \cite{Dereventsov2020} embeds a number of changes with respect to the previous literature. 

At first, as it was done in \cite{zhang2020novel}, it replaces Gaussian Smoothing with Directional Gaussian Smoothing. Where, instead of doing a single $d-$dimensional Gaussian Smoothing, $d$ $1-$dimensional Gaussian Smoothing are performed, which can be approximated via Gauss-Hermite quadrature. 

To achieve this, a function $G:\mathbb{R} \rightarrow \mathbb{R}$ is defined as a cross section of $F$ along $\xi$ as

\begin{equation}
	G(y \mid x, \xi)=F(x+y \xi), \quad y \in \mathbb{R} 
\end{equation}

It is further defined the Gaussian smoothing of $G(y)$, denoted as $G_\sigma (y)$, by

\begin{equation}
	G_{\sigma}(y \mid x, \xi):=\frac{1}{\sqrt{2 \pi}} \int_{\mathbb{R}} G(y+\sigma v \mid x, \xi) \mathrm{e}^{-\frac{v^{2}}{2}} d v=\mathbb{E}_{v \sim \mathcal{N}(0,1)}[G(y+\sigma v \mid x, \xi)] 
\end{equation}

which is the Gaussian smoothing of $F(x)$ along $\xi$ near $x$. Using a one-dimensional expectation, the derivative of $G(y \mid x, \xi)$ at $y=0$, is 

\begin{equation}\label{DerivativeSmoothingSection}
	\mathcal{D}\left[G_{\sigma}(0 \mid x, \xi)\right]=\frac{1}{\sigma} \mathbb{E}_{v \sim \mathcal{N}(0,1)}[G(\sigma v \mid x, \xi) v] 
\end{equation}

Now let $\Xi = \left(\xi_1,\dots,\xi_d\right)$ an orthonormal basis, than the following vector Directional Gaussian Smoothing gradient (henceforth DGS gradient), can be defined:

\begin{defn}[DGS Gradient]\label{DGSGradient}

	\begin{equation}
		\nabla_{\sigma, \xi}[F](x):=\left[\mathcal{D}\left[G_{\sigma}\left(0 \mid x, \xi_{1}\right)\right], \cdots, \mathcal{D}\left[G_{\sigma}\left(0 \mid x, \xi_{d}\right)\right]\right] \Xi 
	\end{equation}

\end{defn}

At this point, the key is that each of the components of the DGS gradient requires a one-dimensional integral, which is approximated with high accuracy with the Gauss-Hermite rule. Hence we obtain the following estimator for $\mathcal{D}\left[G_{\sigma}(0 \mid x, \xi)\right]$:

\begin{equation}\label{GHApproximationSmothedSection}
	\tilde{\mathcal{D}}^{M}\left[G_{\sigma}(0 \mid x, \xi)\right]=\frac{1}{\sqrt{\pi} \sigma} \sum_{m=1}^{M} w_{m} F\left(x+\sqrt{2} \sigma v_{m} \xi\right) \sqrt{2} v_{m} 
\end{equation}

Thus we define the DGS Gradient Estimator as 

\begin{defn}[DGS Gradient Estimator]\label{DGSGradientEstimator}

	\begin{equation}\label{DGSGradientEstimatorEq}
		\widetilde{\nabla}_{\sigma, \Xi}^{M}[F](x)=\left[\widetilde{\mathcal{D}}^{M}\left[G_{\sigma}\left(0 \mid x, \xi_{1}\right)\right], \cdots, \widetilde{\mathcal{D}}^{M}\left[G_{\sigma}\left(0 \mid x, \xi_{d}\right)\right]\right] \Xi 
	\end{equation}

\end{defn}

The authors then fix the balance between exploration and exploitation by assigning, at each iteration, the previous estimated
gradient as first vector of the new set of direction, keeping the remaining directions randomly chosen. It then
estimates for each direction a Local Lipschitz Constant to update, with a momentum-like mechanism, the learning rate like in \ref{learning rate}.

\begin{algorithm}[!ht]
	\caption{ASGF algorithm}\label{ASGF}
	
	\hspace*{\algorithmicindent} \textbf{Input:} function $f$, state $x_0$, smoothing parameter $\sigma_0$, termination tolerance $\varepsilon$, hyper-parameters $\alpha, k$, total iterations \textit{maxiter}\\
	\hspace*{\algorithmicindent} \textbf{Output:} point of minimum $\Tilde{x}$
	\begin{algorithmic}[1]
		
		\State Set $\Tilde{x} = x_0$, $\sigma = \sigma_0$ and let $\Xi$ be a random orthonormal basis
		\For{$i=0$ to \textit{maxiter}}
		\For{$j=1$ to $d$}
		\State compute $ \widetilde{\mathcal{D}}^{M}\left[G_{\sigma}\left(0 \mid x, \xi_{j}\right)\right] $ by (\ref{GHApproximationSmothedSection}) and estimate local Lipschitz constants $L_j$
		\EndFor
		\State Average Lipschitz constant $L_\nabla$ and compute learning rate $\lambda$
		\State Assemble the gradient $ \widetilde{\nabla}_{\sigma, \Xi}^{M}[F](x_i) $ by (\ref{DGSGradientEstimatorEq}) and update $ x_{i+1}=x_{i}-\lambda \widetilde{\nabla}_{\sigma, \Xi}^{M}[F](x_i) $
		\State Add the gradient estimate to $\mathbf{G}$
		\If{$f\left( x_{i+1}\right) < f(\Tilde{x})$}
		\State Update $\Tilde{x} = x_{i+1}$
		\EndIf
		\If{$\norm{x_{i+1} - x_i}_2 <\varepsilon$}
		\State \textbf{break}
		\Else
		\State Update smoothness parameter $\sigma$ and the search directions $\Sigma$ by Algorithm (\ref{subroutineASGF})
		\EndIf
		\EndFor
		
	\end{algorithmic}
\end{algorithm}

\begin{algorithm}[!ht]
	\caption{Parameter update}\label{subroutineASGF}
	
	\hspace*{\algorithmicindent} \textbf{Input:} smoothing parameter $\sigma$, gradient $ \widetilde{\nabla}_{\sigma, \Xi}^{M}[F](x_i)$, local Lipschitz constants $L_1,\dots,L_d$\\
	\hspace*{\algorithmicindent} \textbf{Data:} number of resets $r$ and reset factor $\rho$, decay rate $\gamma_\sigma$, threshold parameters $A,B$ and their change rates $A_+, A_-, B_+, B_-$\\
	\hspace*{\algorithmicindent} \textbf{Output:} smoothing parameter $\sigma$ and directions $\Sigma$
	\begin{algorithmic}[1]
		
		\If{$r>0$ \textbf{and} $\sigma < \rho \sigma_0$}
		\State Assign $\Xi$ to be a random orthonormal basis and set $\sigma = \sigma_0$
		\State Set $A,B$ to their initial values and set $r = r-1$
		\Else
		\State Update directions $\mathcal{D}$
		\If{\( \max _{1 \leq j \leq d}\left| \widetilde{\mathcal{D}}^{M}\left[G_{\sigma}\left(0 \mid \boldsymbol{x_i}, \xi_{j}\right)\right] / L_{j}\right|<A \)}
		\State Decrease smoothing $\sigma = \sigma * \gamma_{\sigma}$ and lower threshold $A=A*A_-$
		\ElsIf{\( \max _{1 \leq j \leq d}\left| \widetilde{\mathcal{D}}^{M}\left[G_{\sigma}\left(0 \mid \boldsymbol{x_i}, \xi_{j}\right)\right] / L_{j}\right|> B \)}
		\State Increase smoothing $\sigma = \sigma / \gamma_{\sigma}$ and upper threshold $B=B*B_+$
		\Else
		\State Increase lower threshold $A=A*A_+$ and decrease upper threshold $B=B*B_-$
		\EndIf
		\EndIf
		
	\end{algorithmic}
\end{algorithm}
